% Context from chapters-tex/denoising.tex; for mathematical reference, not compilation.
                                                                                                              

  
\paragraph{Translation invariant wavelet frame.}

Cycle spinning also has a frame interpretation: replace the orthogonal basis $\Bb = \{\psi_m\}$ by its translated copies,
\eql{\label{eq-ti-wavframe}
	\Bb_{\text{inv}} = \{ (\psi_m)_\tau \}_{m, \tau \in \De}.
}
Although $\Bb_{\text{inv}}$ is no longer an orthogonal basis, orthonormality of each translated copy gives the energy and reconstruction identities
\eq{
	\norm{f}^2 = \frac{1}{|\De|} \sum_{m, \tau \in \De} |\dotp{f}{(\psi_m)_\tau}|^2
	\qandq
	f = \frac{1}{|\De|} \sum_{m, \tau \in \De} \dotp{f}{(\psi_m)_\tau} (\psi_m)_\tau.
}
Such redundant families are called tight frames.

One can then define a translation invariant thresholding denoiser
\eql{\label{eq-ti-thresh}
	\Dd_{\text{inv}}(f) = \frac{1}{|\De|} \sum_{m, \tau \in \De} S_T(\dotp{f}{(\psi_m)_\tau}) (\psi_m)_\tau.
}
This estimator equals the cycle spinning estimator defined in \eqref{eq-cycle-spinning}.

Counted with multiplicity, the translated frame has $|\De| |\Bb|$ elements.
For a basis of signals with $N$ samples and a lattice of $|\De|=N$ shifts, this gives up to $N^2$ vectors in $\Bb_{\text{inv}}$. In a hierarchical construction such as a wavelet basis, however, different shifts can produce the same atom:
\eq{
	(\psi_m)_\tau = (\psi_{m'})_{\tau'}
	\qforq
	m \neq m'
	\qandq
	\tau\neq \tau',
}
The number of distinct vectors in $\Bb_{\text{inv}}$ can therefore be much smaller than $N^2$.
For instance, for an orthogonal wavelet basis, one has
\eq{
	(\psi_{j,n})_{k 2^j} = \psi_{j,n+k},
}
so many translated atoms coincide. For a signal of length $N$, an undecimated transform has $\log_2N$ detail arrays of length $N$, plus a coarse array. For an $N_0\times N_0$ image ($N=N_0^2$), it has $3\log_2N_0$ detail arrays of $N$ entries, plus a coarse array. Appropriate multiplicity weights must be retained when replacing the full translated family by distinct atoms.

The fast undecimated wavelet transform, often called the ``\`a trous'' transform, computes these coefficients in $O(N\log N)$ operations. Each detail array is indexed directly by every pixel location:
\eq{
 d_j^\omega[n]=\dotp{f}{(\psi_{j,0}^\omega)_n},\qquad
 n\in\{0,\ldots,N_0-1\}^2,\quad\omega\in\{V,H,D\},
}
where the wavelets and shifts are discrete and appropriately normalized. Each array is a band-pass filtered image, as illustrated in Figure~\ref{fig-tiwav}.

\myfigure{
\tabquatre{
\image{denoising}{.23}{tiwav-2d-image}&
\image{denoising}{.23}{tiwav-2d-j2-h}&
\image{denoising}{.23}{tiwav-2d-j2-v}&
\image{denoising}{.23}{tiwav-2d-j2-d}\\
$f$ & $j=-8, \om=H$ & $j=-8, \om=V$ & $j=-8, \om=D$ \\
&
\image{denoising}{.23}{tiwav-2d-j1-h}&
\image{denoising}{.23}{tiwav-2d-j1-v}&
\image{denoising}{.23}{tiwav-2d-j1-d}\\
& $j=-7, \om=H$ & $j=-7, \om=V$ & $j=-7, \om=D$ 
}
}{ 
	Translation invariant wavelet coefficients.   
}{fig-tiwav}

Figure \ref{fig-tiwav-coefs-thresh} illustrates hard thresholding of these translated coefficients.

\myfigure{
\image{denoising}{.3}{tiwav-coefs}
\image{denoising}{.3}{tiwav-coefs-thresh}
}{ 
	Left: translation invariant wavelet coefficients, for $j=-8, \om=H$, right: thresholded coefficients.  
}{fig-tiwav-coefs-thresh}



                                                     
\subsection{Other Shrinkage Rules}

Other shrinkage rules offer different compromises between retaining large coefficients and suppressing noise. We describe two examples whose performance depends on the coefficient distribution of $f$ in the chosen basis.

  
\paragraph{Semi-soft thresholding.}

Semi-soft thresholding interpolates between hard and soft thresholding through a parameter $\mu>1$:
\eq{
	S_T^\th(x) = g_{\frac{1}{1-\th}}(x)
	\qwhereq
	g_\mu(x) = 
	\choice{
		0 \qifq |x|<T\\
		x \qifq |x|>\mu T\\
		\sign(x)\frac{\mu(|x|-T)}{\mu-1} \text{ otherwise.}
	}
}
For $0<\theta<1$, set $\mu=1/(1-\theta)$. The limits $\theta\downarrow0$ and $\theta\uparrow1$ give hard and soft thresholding, respectively, away from the threshold boundary.
Figure~\ref{fig-semisoft} shows an example of this nonlinearity.

\myfigure{
\image{denoising}{.45}{semisoft-thresholders}
\image{denoising}{.45}{stein-thresholder}
}{ 
	Left: semi-soft thresholder, right: Stein thresholder.   
}{fig-semisoft}

Figure \ref{fig-semisoft-optimal} shows an SNR improvement for a suitable choice of $\mu$, although the visual difference from hard and soft thresholding is small.

\myfigure{
\image{denoising}{.45}{semisoft-optimal-image}
\image{denoising}{.4}{semisoft-optimal-curve}
}{ 
	 Left: SNR as a function of $\mu$ and $T/\si$. Right: SNR versus $\mu$, with $T/\si$ optimized separately for each $\mu$.  
}{fig-semisoft-optimal}

  
\paragraph{Stein thresholding.}

Stein thresholding is defined using a quadratic attenuation of large coefficients
\eq{
 	S_T^{\text{Stein}}(x) = \max\pa{1-\frac{T^2}{|x|^2},0}x.
}
The value at $x=0$ is defined to be zero. This should be compared with the linear attenuation of soft thresholding
\eq{
 	S_T^1(x) = \max\pa{1-\frac{T}{|x|},0}x.
}
For large coefficients, Stein thresholding has the following shrinkage behavior:
\eq{
	|S_T^{\text{Stein}}(x)-x| \rightarrow 0
	\quad\text{ whereas }\quad
	|S_T^1(x)-x| \rightarrow T,
}
as $x \rightarrow \pm \infty$. Its deterministic shrinkage tends to zero, whereas soft thresholding retains a fixed shrinkage. This limiting property does not imply unbiased estimation at finite amplitudes.

\myfigure{
\image{denoising}{.45}{stein-optimal-curve}
}{ 
	 SNR as a function of $T/\si$ for Stein thresholding.  
}{fig-stein-optimal-curve}

Stein and hard thresholding give similar results in the translation-invariant natural-image experiments shown here.

                                                     
\subsection{Block Thresholding}

The nonlinear thresholding methods above are diagonal estimators, since they attenuate each coefficient separately
\eq{
	\tilde f = \sum_m A_T^q(\dotp{f}{\psi_m}) \dotp{f}{\psi_m} \psi_m
} 
where 
\eq{
	A_T^q(x) = 
	\choice{
		\max(1-T^2/|x|^2,0) \qforq q=\text{Stein}\\
		\max(1-T/|x|,0) \qforq q=1\text{ (soft)}\\
		1_{|x|>T} \qforq q=0\text{ (hard)}
	}
}
Define every attenuation factor to be zero at $x=0$ when $T>0$. Block thresholding shares an attenuation factor across a group of coefficients, exploiting their statistical dependence. This is useful for natural images, where edges create clusters of large wavelet coefficients. Grouping also helps suppress isolated noise coefficients in otherwise smooth regions.

Partition the coefficients into disjoint blocks; for example, in a 2-D wavelet transform,
\eq{
	\{ (j,n,\om) \}_{j,n,\om} = \bigcup_k B_k,
}
where each $B_k$ contains $s\times s$ neigh